\section{Динамические игры с полной информацией. Непрерывные стратегии}

В данной лекции рассматриваются динамические игры, в которых игроки принимают решения последовательно, обладая полной информацией о предшествующих ходах, а пространства стратегий могут быть непрерывными.

\subsection{Поведенческие аспекты и отклонения от рациональности}

\begin{example}[Игра «Ультиматум»]
Экспериментатор предлагает Игроку 1 поделить $100$ рублей между собой и Игроком 2. Игра протекает в два этапа:
\begin{enumerate}
    \item Игрок 1 делает предложение о дележе: $p$ рублей --- себе, $100 - p$ рублей --- партнеру.
    \item Игрок 2 либо принимает предложение (тогда выигрыши реализуются согласно предложению), либо отказывается (в этом случае оба игрока получают $0$).
\end{enumerate}
С точки зрения строгой рациональности и поиска совершенного по подыграм равновесия, Игрок 2 должен соглашаться на любую строго положительную сумму ($100 - p > 0$). Предвидя это, Игрок 1 предлагает минимально возможную положительную сумму (или $0$, если пространство непрерывно и безразличие разрешается в пользу принятия). 

Однако эмпирические данные показывают иную картину:
\begin{itemize}
    \item В среднем люди предлагают партнеру значимую долю (от $30$ до $70$ рублей).
    \item Низкие предложения (менее $20-30\%$) чаще всего отвергаются Игроком 2.
\end{itemize}
Причины таких отклонений кроются в факторах, выходящих за рамки максимизации монетарного выигрыша: обостренное чувство справедливости, репутационные издержки, эмоциональные реакции (обида), а также социальные установки (гражданский долг наказать несправедливого партнера).
\end{example}

\begin{example}[Игра «Доверие»]
Рассмотрим игру с непрерывным пространством стратегий:
\begin{itemize}
    \item В момент $t = 1$ Игрок 1 решает, какую сумму $s_1 \in [0, 100]$ передать Игроку 2.
    \item В момент $t = 2$ переданная сумма утраивается. Игрок 2 решает, какую сумму $s_2 \in [0, 3s_1]$ вернуть Игроку 1.
\end{itemize}
Функции выигрыша игроков имеют вид:
\begin{align}
    u_1(s_1, s_2) &= 100 - s_1 + s_2 \\
    u_2(s_1, s_2) &= 3s_1 - s_2
\end{align}
Поиск равновесия осуществляется методом обратной индукции. Оптимальный ответ Игрока 2 на любой ход Игрока 1 заключается в минимизации $s_2$, следовательно, $\check{s}_2(s_1) = 0$. Подставляя это в функцию выигрыша первого игрока, получаем $u_1(s_1, \check{s}_2(s_1)) = 100 - s_1$. Максимизируя данное выражение по $s_1$, Игрок 1 выбирает равновесную стратегию $s_1^* = 0$.
\end{example}

\subsection{Координация и анализ подыгр}

\begin{example}[Модель экспроприации]
В игре участвуют три игрока: Диктатор (Игрок 1) и две группы граждан (Группа A и Группа B). 
\begin{enumerate}
    \item Диктатор принимает решение: какую группу ограбить (A, B) или не грабить никого (N).
    \item Обе группы одновременно решают: восставать (R) или не восставать (N).
\end{enumerate}
Если экспроприации не происходит, выигрыши всех игроков равны $0$. В случае экспроприации возникают две симметричные подыгры. Выигрыши в матрицах указаны в порядке: \textit{(Диктатор, Группа A, Группа B)}.

\begin{table}[h]
\centering
\renewcommand{\arraystretch}{1.3}
\begin{NiceTabular}{cc|cc}
\toprule
& & \multicolumn{2}{c}{\textbf{Группа B}} \\
& & \textit{R} & \textit{N} \\
\midrule
\multirow{2}{*}{\textbf{Группа A}} 
& \textit{R} & $-k, b, b$ & $x, -x - c, 0$ \\
& \textit{N} & $x, -x, -c$ & $x, -x, 0$ \\
\bottomrule
\end{NiceTabular}
\caption{Подыгра 1: Диктатор грабит Группу A}
\end{table}

\begin{table}[h]
\centering
\renewcommand{\arraystretch}{1.3}
\begin{NiceTabular}{cc|cc}
\toprule
& & \multicolumn{2}{c}{\textbf{Группа B}} \\
& & \textit{R} & \textit{N} \\
\midrule
\multirow{2}{*}{\textbf{Группа A}} 
& \textit{R} & $-k, b, b$ & $x, -c, -x$ \\
& \textit{N} & $x, 0, -x - c$ & $x, 0, -x$ \\
\bottomrule
\end{NiceTabular}
\caption{Подыгра 2: Диктатор грабит Группу B}
\end{table}

В каждой из подыгр существует ровно два равновесия Нэша в чистых стратегиях: $(R, R)$ и $(N, N)$. Для нахождения совершенных по подыграм равновесий (SPNE) необходимо рассмотреть все возможные комбинации равновесий в подыграх и определить оптимальный ход Диктатора на первом этапе.

\begin{table}[h]
\centering
\renewcommand{\arraystretch}{1.3}
\begin{NiceTabular}{cccc}
\toprule
\textbf{№} & \textbf{Равновесие в Подыгре A} & \textbf{Равновесие в Подыгре B} & \textbf{Опт. ход Диктатора} \\
\midrule
1 & $(R, R)$ & $(R, R)$ & $N$ \\
2 & $(R, R)$ & $(N, N)$ & $B$ \\
3 & $(N, N)$ & $(R, R)$ & $A$ \\
4 & $(N, N)$ & $(N, N)$ & $A$ или $B$ \\
\bottomrule
\end{NiceTabular}
\caption{Анализ совершенных по подыграм равновесий}
\end{table}

Всего в игре существует $5$ совершенных по подыграм равновесий в чистых стратегиях. Важно отметить, что стратегия граждан «всегда восставать, если меня пытаются ограбить» сама по себе не формирует равновесие без координации действий обеих групп.
\end{example}

\subsection{Политическая экономика: Демократизация «сверху»}

Рассмотрим теоретико-игровую модель политических транзитов, предложенную Д. Аджемоглу и Дж. Робинсоном (2006).

\begin{definition}[Базовые параметры модели]
В обществе существуют два игрока (класса): бедные (доля населения $\lambda \ge \frac{1}{2}$, доход $y^p$) и богатые (доля $1 - \lambda$, доход $y^r > y^p$). Средний доход в обществе составляет:
\begin{equation}
    \bar{y} = \lambda y^p + (1 - \lambda)y^r
\end{equation}
\end{definition}

Рассмотрим сценарий революции. Если бедные устраивают революцию, весь доход за вычетом уничтоженной доли $(1 - \mu)$ достается им. Выигрыши классов после революции составят:
\begin{equation}
    u^p(R) = \frac{\mu\bar{y}}{\lambda}, \quad u^r(R) = 0
\end{equation}

Если революции не происходит, устанавливается ставка налога $\tau \in [0, 1]$. Собранные налоги перераспределяются поровну, однако возникает чистая потеря общества (deadweight loss), равная $\frac{1}{2}\tau^2\bar{y}$. Посленалоговый доход представителя группы $i \in \{r, p\}$ равен:
\begin{equation}
    u^i(N, \tau) = (1 - \tau)y^i + \left(\tau - \frac{1}{2}\tau^2\right)\bar{y}
\end{equation}
Максимизируя данную функцию, получаем предпочитаемые ставки налогов: для бедных $\tau^{p*} = \frac{\bar{y} - y^p}{\bar{y}}$, для богатых $\tau_r^* = 0$. Функция выигрыша богатых $u^r(N, \tau)$ монотонно убывает по $\tau$.

\begin{lemma}[Угроза революции при диктатуре]
Если издержки революции низки и выполняется неравенство $\mu > \bar{\mu}$, где
\begin{equation}
    \bar{\mu} = \frac{\lambda}{2\bar{y}^2} \left( (y^p)^2 + \bar{y}^2 \right)
\end{equation}
то бедные предпочтут революцию при любой ставке налога $\tau \in [0, 1]$, так как $u^p(N, \tau_p^*) < u^p(R)$. Если же $\mu \le \bar{\mu}$, революцию можно предотвратить.
\end{lemma}

В равновесии при диктатуре (когда элиты контролируют ставку $\tau$ и стремятся избежать революции) богатые предложат минимальную ставку $\tau$, при которой $u^p(R) \le u^p(N, \tau)$. Налог будет строго положительным, если выполняется условие $y^p < \frac{1 - \lambda}{\lambda(1 - \mu)}y^r$. Таким образом, даже при диктатуре может наблюдаться перераспределение доходов для удержания бедных от бунта.

\subsubsection{Проблема невыполнимости обещаний (Commitment problem)}

Модель усложняется, если учесть, что обещания диктатора могут быть нарушены после того, как угроза революции минует. Рассмотрим следующую динамику:

\begin{enumerate}
    \item Богатые решают: сохранить политический контроль ($N$) или устроить демократизацию ($D$).
    \item Если контроль сохранен, богатые устанавливают текущую ставку $\tau^{r1}$.
    \item Бедные решают: устраивать революцию ($R$) или нет ($NR$).
    \item Если революции не было, Природа с вероятностью $q$ позволяет богатым пересмотреть налог (в этом случае они установят $\tau^{r2} = 0$), нарушив обещание. С вероятностью $1-q$ налог остается $\tau^{r1}$.
    \item Если на шаге 1 произошла демократизация, власть переходит к бедным, которые устанавливают свою идеальную ставку $\tau^p = \tau^{p*}$.
\end{enumerate}

\begin{figure}[h]
\centering
\begin{istgame}
\xtdistance{15mm}{40mm}
\istroot(0){Богатые (Элита)}
  \istb{D \text{ (Демократизация)}}[al]
  \istb{N \text{ (Сохранение контроля)}}[ar]
  \endist
\xtdistance{15mm}{25mm}
\istroot(1)(0-1){Бедные}
  \istb{R}[al]{(u^r(R), u^p(R))}
  \istb{NR}[ar]{(u^r(N, \tau^{p*}), u^p(N, \tau^{p*}))}
  \endist
\istroot(2)(0-2){Богатые}
  \istb{\tau^{r1}}[r]
  \endist
\istroot(3)(2-1){Бедные}
  \istb{R}[al]{(u^r(R), u^p(R))}
  \istb{NR}[ar]
  \endist
\istroot(4)(3-2){Природа}
  \istb{1-q}[al]{(u^r(N, \tau^{r1}), u^p(N, \tau^{r1}))}
  \istb{q}[ar]{(u^r(N, 0), u^p(N, 0))}
  \endist
\end{istgame}
\caption{Дерево игры: демократизация и проблема доверия}
\end{figure}

\begin{theorem}[Условия демократизации]
Из-за вероятности обмана $q$, богатые не могут достоверно пообещать перераспределение. Чтобы бедные не устроили революцию при сохранении власти богатыми, должна существовать такая ставка $\tau^{r1}$, что ожидаемый выигрыш бедных будет не меньше выигрыша от революции:
\begin{equation}
    (1 - q)u^p(N, \tau^{r1}) + q u^p(N, 0) \ge u^p(R)
\end{equation}
\end{theorem}

Возникают ситуации, когда даже максимальная ставка $\tau^{p*}$ не удовлетворяет этому условию:
\begin{equation}
    (1 - q)u^p(N, \tau^{p*}) + q u^p(N, 0) < u^p(R)
\end{equation}
Это неравенство эквивалентно условию $\mu > \mu_1$, где:
\begin{equation}
    \mu_1 = \frac{\lambda y^p}{\bar{y}} + \lambda(1 - q)\frac{(\bar{y} - y^p)^2}{2\bar{y}^2}
\end{equation}

\textbf{Ключевые выводы:}
\begin{itemize}
    \item Если $u^p(R)$ слишком велико (или $\mu > \bar{\mu}$), происходит \textbf{революция}.
    \item Если издержки революции лежат в диапазоне $\mu \in [\mu_1, \bar{\mu}]$, богатые не могут откупиться налогами из-за проблемы доверия. Единственный способ спастись от революции --- передать политический контроль бедным (\textbf{демократизация}).
    \item Если $\mu \le \mu_1$, элиты остаются у власти, устанавливая компромиссный налог $\tau^{r1*}$ (\textbf{аккомодация}).
\end{itemize}

Невыполнимость обещаний является фундаментальной проблемой в политике. Революцию трудно организовать из-за проблем координации, поэтому у элит всегда существует стимул дать ложные обещания для разрядки обстановки. Осознавая это, общество требует институциональных гарантий --- передачи политической власти.
